\documentclass[11pt]{article}
\usepackage[margin=0.9in]{geometry}
\usepackage{amsmath,amssymb,amsthm,booktabs,mathrsfs}
\usepackage[T1]{fontenc}
\usepackage{lmodern,microtype}
\usepackage[colorlinks=true,linkcolor=blue,urlcolor=blue]{hyperref}
\newcommand{\norm}[1]{\left\lVert#1\right\rVert}
\newtheorem{theorem}{Theorem}[section]
\newtheorem{lemma}[theorem]{Lemma}
\newtheorem{proposition}[theorem]{Proposition}
\newtheorem{corollary}[theorem]{Corollary}
\theoremstyle{remark}
\newtheorem{remark}[theorem]{Remark}
\title{Toward a four-dimensional Yang--Mills construction\\
\large Notebook 36: soft determinants and bounded coarse corrections}
\author{Research working document}
\date{17 October 2026}
\begin{document}
\maketitle
\begin{abstract}
We factor the actual covariant mean-zero determinant of Notebook 35
into a soft scalar/vector pair and two determinants with one color
fiber per final block. Both coarse determinants have uniformly bounded
actual Maxwell coefficients as the block size grows. Their flat gaps
alone would not prove this: we establish rooted physical source bounds,
with all powers of the fine spacing and all fiber dimensions retained.

A new gauge-covariant Hessian estimate controls the actual effective
classical action through its first two background derivatives, uniformly
in preceding depth. Local formal path gauges isolate curvature; an
ordered Laurent factorization retains the shifted lattice components.
This yields the relative vertices needed for the vector correction.
The auxiliary off-shell family is identified with the actual one only
through the proved on-shell affine jets.

The fine soft pair has uniformly bounded short-time and long-time
affine heat contributions. The latter follows from a local cubic
resolvent formula and Schatten bounds that retain the block projector.
Consequently only the intermediate heat interval can supply the
required logarithm.

The resulting reduction retains the full effective action and original
constraint normalization. It does not yet compute the remaining fine
determinant's logarithmic coefficient or establish continuum
Yang--Mills existence.
\end{abstract}

\section*{Conventions and proof boundary}

The primitive odd ratio \(\ell\) is fixed and \(K=\ell^k\).
Current coordinates use counting pairings; their spacing in final
physical units is \(K^{-1}\). Normalization \(\mathcal N_0\) subtracts
the flat value, and \([t^2]\) means a Taylor coefficient.
All actual identities refer to the homogeneous stationary germ
constructed in the preceding notebooks. Exact finite determinant
identities and finite flat spectral estimates are distinguished from
uniform spatial and curvature estimates on a chosen joint infinite
classical structural system. No uniqueness of that system is assumed.

For a trace-normalized Cartan element \(H\), put
\[
 C_H=\sum_{q\in\operatorname{spec}(\operatorname{ad}_H)}q^2,
 \qquad \alpha_g=\frac{C_H}{48\pi^2}.
\]
The target remains the actual coefficient
\(22\alpha_g\log K+O(1)\), uniformly in preceding depth.
The finite-part estimates below advance that comparison; they do not
substitute a narrower objective for the existence problem.

